\subsection{Событийная модель}
\begin{itemize}
  \item \verb|function callback(event)| — обработчик системных событий.
  \item Примеры: \verb|Ev.COPTER_LANDED| (посадка), \verb|Ev.LOW_VOLTAGE2| (низкое напряжение), \verb|Ev.POINT_REACHED|, \verb|Ev.TAKEOFF_COMPLETE|.
  \item Внутри \verb|callback| быстро реагируйте на событие; не используйте \verb|sleep()|.
\end{itemize}
\textbf{Примеры:}
\begin{codefile}{lua/examples/05_leds/leds_events_landed.lua}
\end{codefile}
\textbf{Пояснение.} В этом листинге функция \verb|callback(event)| отслеживает событие \verb|Ev.COPTER_LANDED| и при посадке выключает все светодиоды, проходя по индексам линейки. Основной принцип — простая реакция на одно конкретное событие автопилота.
\begin{codefile}{lua/examples/05_leds/leds_events_low_voltage.lua}
\end{codefile}
\textbf{Пояснение.} Во втором примере создаётся периодический таймер \verb|mainTimer|, а \verb|callback(event)| реагирует на \verb|Ev.LOW_VOLTAGE2|: останавливает таймер и через \verb|Timer.callLater| включает красный цвет на всех светодиодах. Основной принцип — объединить событийную модель и таймеры для безопасной индикации аварийного состояния.
